\medskip\noindent\textbf{Supplementary Note 10. Analysis choices and limitations.}\par\smallskip
\begingroup\small\setlength{\parskip}{3pt}
The study is exploratory. The sweep was planned as a pilot to choose six strengths per penalty family and decoder for a fixed-strength benchmark with five new seeds. That benchmark was not run. Descriptive ranking replaced a pooled reference-critic versus discriminator contrast with a margin and multiplicity correction. The ranking rules were written after the immune, ATAC small and simulation 1 sweep scores had informed strength and input choices, but before sweep scores existed for the other five. Scores of the scvi-tools adversary, scANVI and sysVI on lung, pancreas and simulation 2, and of at least one scvi-tools adversary fit on human--mouse immune, already existed. Extending the sweep used the same strengths, seeds and training settings.

Applying these rules to the sweep required several changes. All fitted strengths replaced the six selected strengths, and three seeds replaced five. A leave-one-seed-out error therefore required two of three readings rather than four of five. The sweep itself supplied $b^*$, the common batch maximum and low-signal flags. A curve starting at or above $b^*$ gave no reading. Transferred strengths had to be fitted in the target task and failure-free in every training task. Baselines retained their own seeds and entered despite one missing setting. Collapse rates used all fitted strengths. Label transfer was computed on every penalty fit, and the values at the matched strength were read from those results.

The input pilots' predefined criteria retained posterior means without standardisation and were not extended beyond their two tasks. The low-signal exclusion was adopted after the pilots. Scanorama's neighbour grid was reduced for computational feasibility before its baseline comparison.

After all sweep scores were read, the share of non-dominated strengths and Pareto-frontier summaries were added. The hypervolume cap and main-text collapse range were set to the highest common batch score after the uncapped areas and collapse rates were examined. Both unrestricted measures remain in Supplementary Table S8. Scaling sensitivities were specified broadly beforehand, but their computations were defined after scores were seen. The direct biology check was designed after sweep and label-transfer results, before its own readouts were computed. After its first results and label-transfer results were read, their strength rule changed from nearest $b^*$ on any crossing to the first-crossing rule. Supplementary Note 12 compares both readings.

The sweep's readings and ranks were fixed as the replication prediction, with its agreement rule, before any replication fit was scored. Agreement required all four within-pair gaps of 0.04 or more to retain their signs and mean Kendall's $\tau_b$ to exceed zero by a one-sided sign-flip test. One outward strength extension and bottom ranks for remaining non-crossings were specified (Supplementary Note 13). Task-level and task-family sign flips were added after replication results. The slope-adjusted sensitivity was named beforehand, but its formula was defined later. Replication mean ranks were also computed after its results.

Several planned analyses were removed after scores were seen. Local-versus-global batch removal did not distinguish penalties within tasks. Good-strength width mainly reflected seed variability, and best-of-$k$ comparisons depended on grid width and tuning. Cross-task mean label-transfer F1 was removed in favour of task-specific comparisons. Planned follow-ups on composition, reference choice, batch count, architecture, KL warm-up, regularisation, optimiser and update steps, and a batch-separability premise check were not run.

Leave-one-seed-out errors rest on three seeds and a non-smooth crossing statistic, with no coverage guarantee. Strength spacing affects interpolation. Collapse intervals assume independence despite shared seed references. A weak penalty lowers $b^*$ for all penalties, and related tasks receive separate weight in mean ranks. Replication addresses new training seeds on the same data, not generalisation to new datasets.
\par\endgroup
